%%=============================================================================
%% DADOS DO TRABALHO -- FONTE ÚNICA DA VERDADE
%%
%% Este é o ÚNICO arquivo onde os dados do seu trabalho devem ser digitados.
%% Tudo o que aparece na capa, folha de rosto, ficha catalográfica, folha de
%% aprovação e nos metadados do PDF é gerado a partir daqui.
%%
%% Campos marcados como OBRIGATÓRIO que ficarem em branco aparecem no PDF
%% como "[[ PREENCHER: campo ]]" e geram aviso na compilação. Ao gerar a
%% versão final, use a opção `entrega' em main.tex: ela transforma o aviso
%% em erro e impede a compilação com campos pendentes.
%%
%% Para reaproveitar estes dados no corpo do texto, use os acessores:
%%   \TituloTrabalho  \NomeAutor      \NomeOrientador   \NomeCoorientador
%%   \NomeCurso       \NomePrograma   \NomeInstituicao  \AreaConcentracao
%%   \LinhaPesquisa   \LocalTrabalho  \CidadeTrabalho   \AnoTrabalho
%%=============================================================================

%% --- Identificação ---------------------------------------------------------
\titulo{Ativos digitais públicos e matriz de referência do ENEM: avaliação técnica de alternativas de classificação automatizada de um acervo educacional audiovisual}  % OBRIGATÓRIO (título provisório)
\subtitulo{}                                      % opcional: \subtitulo{}
\titulotraduzido{Public digital assets and the ENEM reference matrix: technical evaluation of alternatives for the automated classification of an audiovisual educational collection}  % opcional (para o abstract)

\autor{Mairton Damasceno Cunha}                         % OBRIGATÓRIO
\autorficha{CUNHA, Mairton Damasceno}                   % OBRIGATÓRIO (formato invertido)

\orientador{Prof. Dr. Claudiomar Rolim Filho}     % OBRIGATÓRIO
\coorientador{}                                   % opcional

%% --- Vínculo institucional -------------------------------------------------
\curso{Administração Pública}                    % OBRIGATÓRIO em bacharelado/especialização
\programa{Administração Pública}                 % OBRIGATÓRIO em mestrado/doutorado
\area{}                                          % área de concentração (pós-graduação)
\linhapesquisa{}                                 % linha de pesquisa (pós-graduação)

\cidade{Brasília}                        % usada na folha de aprovação
\local{Brasília-DF}                      % usado na capa e folha de rosto
\data{2026}                              % ano de depósito

%% --- Defesa ----------------------------------------------------------------
%% Informe APENAS a data; a cidade vem de \cidade acima.
\datadefesa{data a definir}              % OBRIGATÓRIO -> "Brasília, data a definir."

%% Banca: \membrobanca[Instituição]{Nome}{Função}   (instituição é opcional)
\membrobanca{Prof. Dr. Claudiomar Rolim Filho}{Orientador}
\membrobanca{A definir}{Examinador}
\membrobanca{A definir}{Examinador}

%% --- Ficha catalográfica ---------------------------------------------------
%% A ficha OFICIAL é elaborada pela Biblioteca Ministro Moreira Alves.
%% Solicite por biblioteca@idp.edu.br e transcreva os dados recebidos aqui.
\cutter{}                                % notação Cutter fornecida pela Biblioteca
\numeropaginas{27}                       % OBRIGATÓRIO: total de folhas
\cdd{351}                                % OBRIGATÓRIO: CDD provisória (Administração pública); confirmar com a Biblioteca
\assuntos{1. Ativos digitais públicos. 2. Sistemas de recomendação. 3. Inteligência artificial na educação. 4. ENEM. I. Título.}

%% --- Resumo ----------------------------------------------------------------
\palavraschave{Ativos digitais públicos. Valor público. Sistemas de recomendação. Inteligência artificial na educação. Recursos educacionais abertos. Classificação curricular automatizada.}   % OBRIGATÓRIO
\keywords{Public digital assets. Public value. Recommender systems. Artificial intelligence in education. Open educational resources. Automated curricular classification.}             % OBRIGATÓRIO

%% --- Sobreposições opcionais -----------------------------------------------
%% Normalmente você NÃO precisa mexer daqui para baixo.

%% Texto da natureza do trabalho. Por padrão é gerado a partir da opção da
%% classe (bacharelado/especializacao/mestrado/doutorado).
% \preambulo{...}

%% Tipo do trabalho na ficha catalográfica (padrão: Monografia/Dissertação/Tese)
% \tipotrabalho{Dissertação}

%% Linha inteira da data de defesa, se o seu programa exigir outra redação
% \datadefesacompleta{Aprovada em ... de 2026.}

%% Logotipo da capa
% \logo{figuras/idp-logo.png}
% \semlogo
